% !LW recipe = make-deck
% Level 9 (after the tracks meet): fast joins. Theory: algorithms and
% their cost; what a correct algorithm is. Lean: a goal theorem. Then
% the stretch goals.
% Build: make (rules in ../common/deck.mk).
\DocumentMetadata{
  lang        = en-GB,
  pdfstandard = ua-2,
  tagging     = on,
}
\documentclass{ltx-talk}

\input{preamble}

\title{Level 9: Fast Joins}
\author{The Proofs and Programs Club}
\institute{Royal Holloway, University of London}

\begin{document}

\leveltitle{9}{Fast Joins}{After the tracks meet: leans on FPIL}
  {After \levellink{6}}{Next: \levellink{7}, if you have not done it}

\levelframe{9}{Level 9 at a Glance}{%
  a hash join and a sort-merge join, on arrays.}{%
  that each computes exactly the join of Level 4F.}{%
  program with \code{Array}, and show Lean that loops terminate &
    \rds{fpil-arrays, fpil-arrterm} \\
  prove a fast function equal to a simple one &
    \rds{fpil-tail, fpil-equiv} \\
  use \code{Std.HashMap}, with \code{Hashable} keys & \rds{fpil-hash} \\
  prove equality up to order and duplicates &
    \rds{mil-finset, ml-multiset} \\
}

% ---------------------------------------------------------------- theory
\theorypart{Fast joins, and when they are correct}
\begin{frame}
  \frametitle{Fast Joins Must Compute the Same Table}
  \begin{itemize}
    \item \textbf{The definition}, as a program: for each row of $R$, for
      each row of $S$, keep the pairs whose keys agree. About
      $|R| \cdot |S|$ steps.
    \item \textbf{Hash join}: put the rows of $S$ in a hash map, by key;
      then look up the partners of each row of $R$. About $|R| + |S|$
      steps.
  \end{itemize}
  \begin{sticky}
    ``Future users of large data banks must be protected from having to
    know how the data is organized in the machine.''\hfill --- E. F. Codd,
    1970
  \end{sticky}
  \slideref{E. F. Codd. \textbf{A Relational Model of Data for Large Shared Data Banks}. CACM 13(6), 1970. \rfsep
    \lnk{ahv6} (section 6.1)}
  \theorytag{9}
\end{frame}

\begin{frame}
  \frametitle{A Correct Join Algorithm Represents the Join}
  \defn{Definition 9.1 (Correct join algorithm)}{A function $J$ from two
    lists to a list is a \emph{correct join} if, whenever $l$ represents
    $R$ and $m$ represents $S$ (Definition 1.2), $J(l, m)$ represents
    $R \join_{k_1 = k_2} S$ (Definition 4F.5).}
  \begin{itemize}
    \item The order in which a hash join finds rows does not matter: only
      the table its list represents.
    \item That is Codd's point: how the data is stored and searched is
      free to change; the answer is not.
  \end{itemize}
  \theorytag{9}
\end{frame}

% ---------------------------------------------------------------- Lean
\leanpart{A correct hash join}
\begin{frame*}
  \frametitle{Definition 9.1 in Lean Is One Theorem}
  \leanblock{goal.lean}
  \sketch
  \code{toRelation} gives the table a list represents. Start from your
  \code{joinOn} of Level 4F, proved correct first, as in Level 6.
  \leantag{9}
\end{frame*}

% ---------------------------------------------------------------- beyond
\partframe{Beyond Level 9}{Stretch goals}

\begin{frame}
  \frametitle{What Relational Algebra Can and Cannot Express}
  \begin{itemize}
    \item \textbf{Can}: exactly the domain-independent queries of
      first-order logic over the tables (\emph{Codd's theorem}).
    \item With only $\sigma$, $\pi$, $\join$ and equalities joined by
      $\wedge$: the \textbf{conjunctive queries}, roughly SQL's
      \code{SELECT}--\code{FROM}--\code{WHERE}. Whether one is contained
      in another is decidable (and NP-complete).
    \item Every operator except $-$ is \textbf{monotone}: more rows in
      the input never mean fewer rows in the answer.
    \item \textbf{Cannot}: transitive closure (who can reach whom), or
      counting. SQL adds \code{WITH RECURSIVE} and \code{GROUP BY}.
  \end{itemize}
  \slideref{E. F. Codd. \textbf{Relational Completeness of Data Base Sublanguages}. In Data Base Systems, Prentice Hall, 1972. \rfsep
    A. K. Chandra, P. M. Merlin. \textbf{Optimal Implementation of Conjunctive Queries in Relational Data Bases}. STOC 1977. \rfsep
    A. V. Aho, J. D. Ullman. \textbf{The Universality of Data Retrieval Languages}. POPL 1979. \rfsep
    \lnk{ahv6}}
  \stretchtag
\end{frame}

\begin{frame}
  \frametitle{Beyond Level 9: Bags, NULL, Aggregation, SQL}
  \begin{itemize}
    \item \textbf{Bags}: SQL keeps duplicates. Which laws survive?
    \item \textbf{NULL}: missing values, and three truth values.
    \item \textbf{Aggregation}: \code{GROUP BY}, \code{COUNT},
      \code{SUM}.
    \item \textbf{SQL}: translate a fragment of SQL into your algebra.
    \item \textbf{Codd's theorem}: the algebra and (safe) first-order
      logic express the same queries.
    \item \textbf{Containment}: decide whether a conjunctive query's
      answers are always among another's.
  \end{itemize}
  \slideref{\lnk{fpil-tqex} \rfsep
    A. K. Chandra, P. M. Merlin. \textbf{Optimal Implementation of Conjunctive Queries in Relational Data Bases}. STOC 1977.}
  \stretchtag
\end{frame}

\begin{frame}
  \frametitle{Others Did This in Coq: Read Them After Trying}
  Verified relational databases, in Coq (now Rocq):
  \begin{itemize}
    \item G. Malecha, G. Morrisett, A. Shinnar, R. Wisnesky. \emph{Toward a
      Verified Relational Database Management System}. POPL 2010.
    \item V. Benzaken, \'E. Contejean, S. Dumbrava. \emph{A Coq
      Formalization of the Relational Data Model}. ESOP 2014.
    \item S. Chu, K. Weitz, A. Cheung, D. Suciu. \emph{HoTTSQL: Proving
      Query Rewrites with Univalent SQL Semantics}. PLDI 2017.
    \item V. Benzaken, \'E. Contejean. \emph{A Coq Mechanised Formal
      Semantics for Realistic SQL Queries}. CPP 2019.
  \end{itemize}
  For the theory of databases: \lnk{ahv}. Addison-Wesley, 1995.
\end{frame}

\end{document}
